\documentclass[10pt, letterpaper]{article}
\usepackage{qrcode}
\usepackage{atveryend}
\usepackage{reddy_cover}
\usepackage{hyperref}

\begin{document}

\vspace{\headerSpacing}
\section*{COVER LETTER}

\vspace{\headerSpacing}
\begin{onecolentry}
	Aadarsha Gopala Reddy\\
	St. Louis, MO 63112\\
	(740) 802-1776\\
	agr@agreddy.com\\
\end{onecolentry}

\vspace{\headerSpacing}
\begin{onecolentry}
	July 24, 2026\\
\end{onecolentry}

\vspace{\headerSpacing}
\begin{onecolentry}
	Hiring Committee\\
	Research Data Analyst\\
	Mike Wiegers Department of Electrical and Computer Engineering\\
	Kansas State University, Manhattan, KS\\
\end{onecolentry}

\vspace{\headerSpacing}
\begin{onecolentry}
	\textbf{Subject:} Application for Research Data Analyst - Mike Wiegers Department of Electrical and Computer Engineering
\end{onecolentry}

\vspace{\headerSpacing}
\begin{onecolentry}
	Dear Hiring Committee,
\end{onecolentry}

\vspace{\headerSpacing}
\begin{onecolentry}
	I am writing to apply for the Research Data Analyst position supporting AI applications for power systems research. I completed my M.S. in Computer Science at Washington University in St.~Louis this past May, and through that program I built hands-on experience across data pipelines, simulation benchmarking, and LLM-based tooling, skills that map directly onto the software development, simulation setup, and research deliverable work this role asks for.
\end{onecolentry}

\vspace{\headerSpacing}
\begin{onecolentry}
	For a course project I implemented DDQN, PPO, and SAC reinforcement learning agents benchmarked against EnergyPlus building-energy simulations, improving efficiency by 35.8\% over baseline controllers, direct experience applying simulation-driven optimization and systematic benchmarking to energy systems. Separately, I built an automated red-blue teaming pipeline using Claude Code for agentic coding workflows, generating visual jailbreak attacks against vision-language models and auto-generating deterministic defenses, mirroring the attack-and-defense structure used in AI-driven power-grid cybersecurity research.
\end{onecolentry}

\vspace{\headerSpacing}
\begin{onecolentry}
	As a Graduate Teaching Assistant for a data engineering course, I taught Snowflake, Airflow, Spark, Kafka, and Flink to a graduate cohort and debugged student pipeline failures weekly, giving me fluency in moving raw data into research-ready form and reviewing others' code. Earlier, as a Data Analyst Intern at Lab714, I wrote SQL queries and built Python/Pandas dashboards over IoT sensor telemetry to surface anomalies and trends for client reporting. What draws me to this role specifically is the chance to apply that combination of data engineering, simulation, and AI tooling to a domain where getting the benchmarking right has concrete, measurable stakes.
\end{onecolentry}

\vspace{\headerSpacing}
\begin{onecolentry}
	I would welcome the chance to discuss this role in more detail. Thank you for considering my application; I look forward to the possibility of contributing to K-State ECE's power systems research.
\end{onecolentry}

\vspace{\headerSpacing}
\begin{onecolentry}
	Sincerely,
\end{onecolentry}

\vspace{\headerSpacing}
\begin{onecolentry}
	Aadarsha Gopala Reddy\\
\end{onecolentry}
\end{document}
